\documentclass[11pt]{article}
\usepackage{ochamath}
\subjectcolor{2E7D5B}
\setsheet{微分積分}{テイラー展開と線形近似　演習 1}{https://university-math-crj.pages.dev/calculus/taylor-linearization/}
\namelinetrue
\begin{document}
\makesheethead{問題}

\problem[マクローリン展開]
次の関数を $x^3$ の項までマクローリン展開せよ。
\[ \text{(1)}\ e^{2x} \qquad \text{(2)}\ \log(1+x) \qquad \text{(3)}\ \frac{1}{1-x} \]
\workspace{50mm}

\problem[線形近似で計算する]
$f(x) = (1+x)^{-1/2}$ の $x=0$ における線形近似を求め、それを使って $\dfrac{1}{\sqrt{1.02}}$ の近似値を小数第2位まで求めよ。
\workspace{45mm}

\problem[ダイオードの小信号抵抗]
ダイオードの電流を $I = I_S e^{V/V_T}$、$V_T = 26\ \mathrm{mV}$ とする。
\begin{enumerate}[label=(\arabic*)]
\item 動作点の電流 $I_D = 2\ \mathrm{mA}$ における小信号抵抗 $r_d$ を求めよ。
\item 動作点から電圧を $1\ \mathrm{mV}$ 増やしたとき、電流の増加量を線形近似で見積もれ。
\end{enumerate}
\workspace{50mm}

\problem[テイラー展開で極限を求める]
次の極限を、テイラー展開を使って求めよ。
\[ \text{(1)}\ \lim_{x\to0}\frac{1-\cos x}{x^2} \qquad \text{(2)}\ \lim_{x\to0}\frac{e^x - 1 - x}{x^2} \]
\workspace{40mm}
\end{document}
